\chapter{Background}
\label{ch:background}

This chapter provides the background knowledge needed to understand the proposed
work and discusses existing applications and research related to propaganda,
misinformation, and harmful-content detection. It also identifies the main research
gaps that motivate the proposed system.

\section{Preliminaries}
Social media has become an important platform for sharing news, opinions, images,
and other information. At the same time, it can also be used to spread propaganda,
misinformation, hate speech, manipulation, and other harmful content. Because
harmful information can spread very quickly, automatic detection systems are
becoming increasingly important.

\subsection{Propaganda}
Propaganda is information designed to influence people's opinions or behavior, often
by using selective information, emotional language, misleading presentation, or
other persuasive techniques. In social media, propaganda can appear as text,
images, videos, hashtags, or coordinated posts.

\subsection{Misinformation and Harmful Content}
Misinformation refers to incorrect or misleading information. Harmful online
content is a broader term that can include hate speech, offensive language,
harassment, manipulated information, and other content that may cause social harm.
These types of content can be difficult to detect automatically because their
meaning often depends on context.

\subsection{Machine Learning for Content Detection}
Machine learning (ML) allows a computer system to learn patterns from previously
labeled examples. Traditional ML methods such as Support Vector Machines (SVM) have
been widely used for text classification. More recent systems use deep learning and
transformer-based language models to understand more complex language
patterns~\cite{duridi2025,plikynas2025}.

\subsection{Deep Learning and Transformer Models}
Deep learning (DL) models can automatically learn useful features from data.
Transformer-based models are particularly useful for language tasks because they
can capture relationships between words over a long text. Models such as BERT and
language models adapted to specific languages have therefore become common in
harmful-content detection~\cite{duridi2025,lora2023,hasan2026}.

\subsection{Graph-Based Detection}
Graph-based approaches represent information as connected nodes and relationships.
They are useful when the relationship between words, users, posts, or accounts is
important. For example, propaganda can depend on relationships between different
parts of a text rather than on individual words alone. Ahmad et al.\ used a
hierarchical graph-based model to capture such relationships in propaganda
detection~\cite{ahmad2025}.

\subsection{Multimodal Detection}
Social media posts often contain more than one type of information. A post may
contain text and an image together. Multimodal detection combines these different
sources of information to make a prediction. Research has shown that image
information can provide important clues for fake-news and hateful-meme
detection~\cite{lin2024,hee2022}.

\subsection{Behavioral and Network Signals}
The behavior of a social-media account can also provide useful information.
Examples include how often an account posts, how it interacts with other users, and
whether it participates in coordinated activity. These signals can be useful when
the actual content is difficult to trust or is deliberately changed to avoid
detection~\cite{declerck2024,sela2025,schneider2026}.

\subsection{Low-Resource Languages}
A low-resource language is a language for which relatively few labeled datasets,
language models, or other computational resources are available. Bengali is an
example where researchers have identified limitations in annotated datasets and
language resources~\cite{lora2023,hasan2026,romim2022}. This makes harmful-content
detection more challenging than in high-resource languages such as English.

\section{Literature Review}
The reviewed literature focuses on computational methods for detecting propaganda,
misinformation, manipulation, hate speech, sarcasm, and related harmful content.
The studies cover traditional machine learning, deep learning, graph models,
transfer learning, multimodal methods, behavioral analysis, and large language
models (LLMs)~\cite{duridi2025,plikynas2025,ahmad2025,schneider2026}.

The literature can be divided into two areas: \textbf{similar applications} and
\textbf{related academic research}.

\subsection{Similar Applications}
Several existing tools and systems address parts of the problem of detecting
misinformation, manipulated information, or suspicious online activity. These
systems show that automated analysis can be provided through web-based interfaces
and APIs, although they do not necessarily perform exactly the same task as the
proposed system.
\begin{itemize}
  \item \textbf{Bot-detection and Coordination Analysis Systems:} Some systems
        focus on identifying automated or coordinated social-media accounts.
        Research-based systems such as those discussed by De~Clerck et al.\ analyze
        account behavior and coordination patterns to identify bot-like
        activity~\cite{declerck2024}. These systems are related to the proposed work
        because they show how user behavior can be used in addition to the content
        of a post.
  \item \textbf{Fact-checking Systems:} Fact-checking platforms and services allow
        users or applications to check whether claims have already been reviewed.
        Google provides Fact Check Tools that support searching and working with
        fact-check information through an API~\cite{google_factcheck}. Google also
        supports the \texttt{ClaimReview} format for representing fact-check
        results in a structured way~\cite{google_claimreview}. Such systems are
        related because they support verification of questionable information,
        although fact-checking is different from automatically detecting
        propaganda techniques.
  \item \textbf{Social-media Misinformation Analysis:} Other applications and
        research systems analyze how information spreads through social networks.
        The work of Sela et al.\ shows that repetition patterns and network
        structures can provide signals for identifying political influence in
        information spread~\cite{sela2025}. These systems are relevant because they
        demonstrate that harmful information can be analyzed at both the content
        level and the network level.
  \item \textbf{Limitations of Existing Applications:} A common limitation of
        existing applications is that they often focus on one specific problem,
        such as fact checking, bot detection, or misinformation analysis. A system
        that combines these ideas with low-resource language and multimodal
        analysis can address a broader set of challenges.
\end{itemize}

\subsection{Related Research}

\subsubsection*{Propaganda Detection Using Machine Learning and Deep Learning}
Traditional machine learning remains useful for some propaganda detection tasks.
Duridi et al.\ found that SVM performed better than AraBERT for propaganda
detection in their Arabic social-media dataset, while AraBERT performed better for
bias detection~\cite{duridi2025}. This result shows that the most advanced model is
not always the best model for every task.

Ahmad et al.\ proposed a Hierarchical Graph-based Integration Network (H-GIN) for
propaganda detection. The model uses graph-based information to capture long-range
and non-adjacent relationships that can be missed by standard sequence models. The
study reported an accuracy of 82\% for its evaluated task~\cite{ahmad2025}.

These studies show that both traditional and advanced models can be useful. The
choice of model depends on the task, the dataset, and the type of information that
needs to be captured.

\subsubsection*{Behavioral and Network-Based Detection}
Recent research has moved beyond text-only detection. De~Clerck et al.\ studied
coordinated and bot-like behavior on Twitter and showed that patterns of account
activity can help identify suspicious behavior~\cite{declerck2024}.

Sela et al.\ also found that repetition patterns and network structure can provide
useful signals for identifying political influence in information
spread~\cite{sela2025}.

Schneider et al.\ further developed this idea by focusing on user behavior rather
than only the content of posts. In their Reddit study, behavioral-policy
representations achieved a macro-F1 of 94.9\%, compared with 91.2\% for text
embeddings in the evaluated actor-detection task~\cite{schneider2026}.

These studies suggest that user behavior and network structure can provide
information that cannot be obtained from text alone.

\subsubsection*{Low-Resource Language Detection}
Low-resource languages face important challenges because of limited data and
language resources. Haider et al.\ studied manipulation detection in low-resource
languages and showed that transfer learning can help when only a small amount of
training data is available in the target language~\cite{haider2020}.

For Bengali, Lora et al.\ developed a transformer-based model for sarcasm
detection. Their results differed across datasets, showing that model performance
can depend strongly on the characteristics of the available corpus~\cite{lora2023}.

Romim et al.\ introduced the BD-SHS benchmark dataset for Bangla hate-speech
detection. The dataset contains manually annotated comments from different social
contexts, and the study reported a best F1-score of 91.0\% for its evaluated
system~\cite{romim2022}.

More recent work by Hasan et al.\ introduced BanglaMultiHate, which considers
several aspects of hate speech, including type, severity, and target. The study
evaluated classical machine learning, BanglaBERT, and LLM-based
approaches~\cite{hasan2026}.

Together, these studies show that successful moderation in Bangla requires larger
datasets, better annotation, and models that understand local language and social
context.

\subsubsection*{Multimodal Detection}
Lin et al.\ proposed a text-image fusion model for fake-news detection. Their
results demonstrate that combining text and image information can improve detection
compared with using only one modality~\cite{lin2024}.

Hee et al.\ studied multimodal hateful memes and found that images can provide
important information for classification. However, they also observed that
multimodal systems may learn unwanted biases and produce false
positives~\cite{hee2022}.

Wang et al.\ investigated image-based political propaganda and showed that visual
elements such as colors can be used in different ways by political
groups~\cite{wang2022}.

These findings suggest that text-only systems may miss important visual
information. Therefore, multimodal analysis is an important direction for future
harmful-content detection.

\subsubsection*{Large Language Models}
Large language models are a newer research direction in propaganda detection.
Pi\~na-Garc\'ia used few-shot LLaMA~3.2 to classify propagandistic tweets and
examine linguistic and coordination patterns in Mexican election-related
data~\cite{pina2025}.

Gaeta et al.\ proposed an LLM-based system that detects propaganda and identifies
the specific parts of a text that contain propagandistic content~\cite{gaeta2025}.

However, LLM performance depends strongly on data quality. Maarouf et al.\ showed
that models trained with weak labels achieved an AUC of 64.03, whereas models
trained with high-quality human annotations achieved an AUC of 92.25. Prompt-based
learning with a smaller set of high-quality samples achieved an AUC of
80.27~\cite{maarouf2024}.

This indicates that better models alone are not enough. High-quality training and
evaluation data are also essential.

\subsection{Summary of Literature}
\label{sec:lit-summary}
Table~\ref{tab:litsummary} provides a systematic comparative summary of the
reviewed literature, outlining the core computational methods, benchmark datasets,
key findings, and critical limitations of prior work.

\begingroup
\small
\setlength{\tabcolsep}{4pt}
\renewcommand{\arraystretch}{1.25}
\begin{longtable}{@{}L{3.1cm}C{0.8cm}L{2.2cm}L{3.5cm}L{3.3cm}@{}}
\caption{Summary of related research papers.}\label{tab:litsummary}\\
\toprule
\textbf{Paper Title \& Citation} & \textbf{Year} & \textbf{Dataset Used} & \textbf{Key Methodology \& Findings} & \textbf{Critical Limitations}\\
\midrule
\endfirsthead
\multicolumn{5}{@{}l}{\small\itshape Table~\thetable\ (continued)}\\
\toprule
\textbf{Paper Title \& Citation} & \textbf{Year} & \textbf{Dataset Used} & \textbf{Key Methodology \& Findings} & \textbf{Critical Limitations}\\
\midrule
\endhead
\midrule
\multicolumn{5}{r@{}}{\small\itshape Continued on next page}\\
\endfoot
\bottomrule
\endlastfoot

\textbf{Spread of propaganda by coordinated communities on social media}~\cite{hristakieva2022} & 2022 & Twitter (Italian elections \& COVID-19) &
Analyzed coordinated link-sharing networks and bot behavior using network graphs. Found that propaganda spreaders exhibit significantly higher coordination than benign users. &
Evaluated only on high-resource languages (Italian/English); text semantics were secondary to network topology; lacks multimodal analysis.\\
\addlinespace
\textbf{BanFakeNews: A dataset for detecting fake news in Bangla}~\cite{hossain2020} & 2020 & BanFakeNews (50,000 Bangla news articles) &
Landmark benchmark Bangla dataset for fake news detection; evaluated SVM, Random Forest, and LSTM architectures. &
Focuses strictly on factual veracity (fake vs.\ real) rather than rhetorical propaganda techniques; purely text-based without social context or coordination tracking.\\
\addlinespace
\textbf{BanMANI: Manipulated social media news in Bangla}~\cite{kamruzzaman2023} & 2023 & BanMANI (social media news posts) &
Categorizes social media news manipulation into multiple classes; demonstrates the efficacy of transfer learning. &
Small dataset size; does not analyze coordinated amplification campaigns or image--text multimodal relationships.\\
\addlinespace
\textbf{Coordinated information campaigns on social media}~\cite{ng2023} & 2023 & Twitter / Reddit multi-domain data &
Multi-faceted framework fusing semantic similarity, temporal proximity, and network clustering to detect coordinated operations. &
Computationally intensive; designed primarily for English; relies on extensive historical user activity graphs often unavailable in low-resource contexts.\\
\addlinespace
\textbf{Propaganda and bias in social media: Israel--Gaza war}~\cite{duridi2025} & 2025 & Arabic X (Twitter) posts &
Evaluated classical ML (SVM) against deep learning (AraBERT); SVM achieved superior performance on propaganda classification while AraBERT excelled at bias detection. &
Restricted to short text posts; does not address multimodal memes or cross-platform campaign coordination.\\
\addlinespace
\textbf{Hierarchical graph-based integration network for propaganda detection}~\cite{ahmad2025} & 2025 & SemEval propaganda corpus (English) &
H-GIN model capturing long-range, non-adjacent token relationships via graph neural networks; reported 82\% accuracy. &
Operates solely on English long-form news articles; high graph construction latency; no support for low-resource or code-mixed text.\\
\addlinespace
\textbf{Text-image multimodal fusion for fake news detection}~\cite{lin2024} & 2024 & Weibo \& Twitter multimodal benchmarks &
Cross-modal attention network fusing visual features (ResNet/ViT) with textual embeddings (BERT); showed multimodal fusion outperforms unimodal baselines. &
Evaluated on Chinese and English benchmarks; vulnerable to false positives when meme text and images are intentionally incongruent.\\
\addlinespace
\textbf{Explaining multimodal hateful meme detection models}~\cite{hee2022} & 2022 & Facebook Hateful Memes dataset &
Investigates multimodal interpretability; reveals that multimodal models frequently latch onto background image biases and surface textual heuristics. &
Demonstrates explainability challenges; lacks application to political propaganda in South Asian contexts.\\
\addlinespace
\textbf{Beyond content: Behavioral policies reveal actors in information operations}~\cite{schneider2026} & 2026 & Twitter \& Reddit coordinated influence sets &
Models user behavioral policies (posting timing, activity patterns) reaching 94.9\% macro-F1, outperforming text embeddings (91.2\%). &
Requires rich temporal account metadata; unviable for anonymous or newly created ephemeral accounts.\\
\addlinespace
\textbf{BD-SHS: Benchmark dataset for online Bangla hate speech}~\cite{romim2022} & 2022 & BD-SHS (Bangla social comments) &
Diverse benchmark across multiple social categories; achieved 91.0\% F1-score using transformer fine-tuning. &
Focuses on abusive language rather than persuasive propaganda techniques; lacks coordination and image features.\\
\addlinespace
\textbf{LLM-based multi-task Bangla hate speech detection}~\cite{hasan2026} & 2026 & BanglaMultiHate dataset &
Multi-task framework evaluating classical ML, BanglaBERT, and few-shot LLMs across hate speech type, severity, and target. &
Shows that while LLMs are promising, domain-specialized models (BanglaBERT) often yield superior and more consistent F1 on local cultural nuances.\\
\addlinespace
\textbf{HQP: A human-annotated dataset for detecting online propaganda}~\cite{maarouf2024} & 2024 & HQP (English articles) &
Proves data quality impact: weak labels gave 64.03 AUC, whereas human-annotated data achieved 92.25 AUC. Prompt-learning with small verified sets achieved 80.27 AUC. &
English-only; underscores the vital necessity of human verification over pure LLM automated annotation.\\
\end{longtable}
\endgroup

\section{Gap Analysis}
The reviewed literature shows several important gaps that can be addressed by the
proposed work:

\subsection{Gap 1: Limited Support for Low-Resource Languages}
Many existing systems are developed and tested mainly on high-resource languages.
Research on Bengali and Bangla shows that limited annotated datasets and language
resources remain major problems~\cite{lora2023,hasan2026,romim2022}. More work is
needed on systems that can effectively handle low-resource languages.

\subsection{Gap 2: Heavy Dependence on Text}
Many detection systems focus mainly on textual information. However, social-media
content often includes images, and important information may appear in the
relationship between text and images~\cite{lin2024,hee2022}. This creates a need
for systems that can analyze multiple modalities together.

\subsection{Gap 3: Limited Use of Behavioral Information}
Traditional content-based systems may fail when users intentionally change or
disguise their language. Behavioral and network-based research shows that user
activity and coordination can provide additional
evidence~\cite{declerck2024,sela2025,schneider2026}. However, these signals are not
always integrated with text and multimodal models.

\subsection{Gap 4: Data Quality and Annotation Problems}
The quality of the training data has a strong effect on model performance. HQP
demonstrates that high-quality human annotation can produce much better results
than weak labeling~\cite{maarouf2024}. Bangla research also emphasizes the
importance of culturally grounded and socially diverse
datasets~\cite{hasan2026,romim2022}.

\subsection{Gap 5: Generalization Across Datasets and Domains}
A model can perform well on one dataset but perform differently on another dataset.
The reviewed literature therefore indicates the need for evaluation across multiple
datasets, domains, and languages rather than relying on a single
benchmark~\cite{ahmad2025}.

\subsection{Gap 6: Explainability}
Many machine-learning and deep-learning systems provide a final classification
without clearly explaining why a particular post was detected as harmful. This can
be especially problematic for propaganda, sarcasm, memes, and culturally sensitive
content. More interpretable systems are therefore
needed~\cite{ahmad2025,hee2022,maarouf2024}.

\subsection{Proposed Research Opportunity}
Based on these gaps, there is an opportunity to develop a system that combines
multiple types of evidence. Such a system could use text analysis together with
image information and, where available, behavioral or network signals. For a
low-resource language such as Bangla, the system could also use language-specific
pretrained models and carefully annotated local datasets.

The main focus should therefore be on building a detection approach that is not
dependent on a single signal or a single dataset. It should be evaluated for
accuracy, robustness, and generalization across different types of online content.

\section{Summary}
This chapter presented the background concepts needed to understand propaganda and
harmful-content detection. It discussed machine learning, deep learning,
graph-based approaches, multimodal analysis, behavioral signals, and the challenges
of low-resource languages.

The literature review showed that researchers are moving from text-only systems
toward approaches that combine content, user behavior, network structure, images,
and large language models~\cite{duridi2025,plikynas2025,lora2023,hasan2026,ahmad2025,lin2024,hee2022,declerck2024,sela2025,schneider2026,romim2022,haider2020,wang2022,pina2025,gaeta2025,maarouf2024}.
The studies also show that no single method works equally well for every task.
Model performance depends on the problem, the dataset, the available language
resources, and the quality of annotation.

The gap analysis identified several areas that still need improvement, especially
low-resource language support, multimodal analysis, behavioral information, data
quality, cross-dataset evaluation, and explainability. These gaps provide the
motivation for developing a more integrated detection system that can work with
low-resource and multimodal social-media content.

% =====================================================================
%                            CHAPTER 3
% =====================================================================